\subsection{Geometric realization}

Consider the functor $\abs{\Delta^\bul}\colon\Delta\to\Top$ which maps $[n]$ to $\abs{\Delta^n}$, and an order preserving map $\phi\colon[n]\to[m]$ to the continuous map
$\abs{\Delta^n}\to\abs{\Delta^m}$ given by
$$ (x_0,\dots,x_n)\in\abs{\Delta^n} \mapsto (y_0,\dots,y_m)\in\abs{\Delta^m},\qquad y_j = \sum_{\phi(i)=j}x_i . $$
For example, the face maps $\delta^i$ define maps $\abs{\Delta^{n-1}}\to\abs{\Delta^n}$ by
$$ \delta^i(x_0,\dots,x_{n-1}) = (x_0,\dots,0,\dots,x_{n-1}), $$
and the degeneracy maps $\sigma^i\colon\abs{\Delta^n}\to\abs{\Delta^{n-1}}$ give
$$ \sigma^i(x_0,\dots,x_n) = (x_0,\dots,x_i+x_{i+1},\dots,x_n) . $$

This then defines a functor
$$ \hom_\Top(\abs{\Delta^\bul},\bul) \colon \Delta^\op\times\Top \to \Set . $$
In particular for every topological space $X$, we have a simplicial set
$$ \S(X) = \hom_\Top(\abs{\Delta^\bul},X) , $$
this is called the {\emphcolor singular set} of $X$.
The face and degeneracy maps of the singular set are then
$$ d_i(\rho) = \rho\circ\delta^i,\qquad s_i(\rho) = \rho\circ\sigma^i . $$

\bnote

    The singular set of $X$ plays a central role in defining the singular homology of $X$.
    Recall that $X$'s chain complex is defined at each coordinate to be the free Abelian group generated by $\S(X)_n$: $C_n(X)=Z[\S(X)_n]$.
    The differentials are then defined in terms of the face maps: $\partial\colon C_n(X)\to C_{n-1}(X)$ is defined on a singular $n$-simplex $\sigma\in\S(X)_n$ by
    $$ \partial\sigma = \sum_{i=0}^n(-1)^id_i\sigma , $$
    and $\partial$ is extended to all of $C_n(X)$ by linearity.

\enote

Geometrically, we view a simplicial set $X$ as a collection of (possibly degenerate) $n$-simplices $X_n$ for all $n<\omega$.
If we attach the geometric $n$-simplex $\abs{\Delta^n}$ to each $n$-simplex $X_n$ we obtain $X_n\times\abs{\Delta^n}$.
We give each $X_n$ the discrete topology, so this becomes a topological space.
The geometric realization of $X$ as a whole should be the union of all these simplices, but this is not enough.
For one, degenerate simplices should not be counted; and we want to glue the faces of simplices together.

So we define the geometric realization as follows.

\bdefn

    For a simplicial set $X$, we define its {\emphcolor geometric realization} to be the space
    $$ \abs X = \coprod_{n<\omega}X_n\times\abs{\Delta^n}/{\sim} $$
    where each $X_n$ has the discrete topology and $\sim$ is the equivalence generated by the relations
    $$ (x,\delta^ip) \sim (d_ix,p),\qquad (x,\sigma^ip) \sim (s_ix,p) . $$

    Here $\delta^i\colon\abs{\Delta^{n-1}}\to\abs{\Delta^n}$ and $\sigma^i\colon\abs{\Delta^n}\to\abs{\Delta^{n-1}}$ are the face and degeneracy maps respectively
    of the standard simplices.

\edefn

For $x\in X_n$, we denote by $(x,\abs{\Delta^n})$ the image of $\set x\times\abs{\Delta^n}$ in $\abs X$.

\blemm

    If $X$ is a simplicial set, then every degenerate simplex of $X$ is the degeneracy of a unique non-degenerate simplex.
    That is, if $x\in X$ is degenerate then there is a unique non-degenerate $y\in X$ such that $x=s_{i_1}\cdots s_{i_k}y$ for some $i_j$s.

\elemm

\bproof

    Clearly if $x$ is degenerate there must exist {\it some\/} $y$ which satisfies this.
    Now suppose $x=Sy=S'y'$ for $y,y'$ non-degenerate and $S,S'$ a composition of degeneracy maps.
    Suppose $S=s_{i_1}\cdots s_{i_k}$, then define $D=d_{i_k}\cdots d_{i_1}$.
    This is a left inverse for $S$, and thus $y=DS'y'$.

    Degeneracy and face maps commute according to the identities laid out in \gotopasteqn{simpid}.
    Thus $y=DS'y'=\tilde S\tilde Dy'$ for some composition of degeneracy maps and face maps.
    Since $y$ is by assumption non-degenerate, $\tilde S$ must be empty so that $y=\tilde Dy'$.

    By symmetry, $y'=\bar Dy$ for some composition of face maps as well.
    Since face maps lower the dimension, this can only happen if $\bar D,\tilde D$ are empty.
    Thus $y=y'$.
    \qed

\eproof

Similarly we have the following result.

\blemm

    Every point $v\in\abs{\Delta^n}$ can be expressed uniquely as $v=\delta^{i_1}\cdots\delta^{i_k}u$ where $u$ is an interior point (all of its coordinates are positive).

\elemm

Call a point $(x,v)\in\abs X$ {\emphcolor non-degenerate} if $x$ is non-degenerate and $v$ is interior.

\blemm

    For a simplicial set $X$, every point in its geometric realization $(x,v)\in\abs X$ is equivalent to a unique non-degenerate point of $\abs X$.

\elemm

\bproof

    Define $\lambda\colon\abs X\to\abs X$ as follows.
    For $x\in X_n$ take $i_1,\dots,i_k$ and $y$ non-degenerate such that $x=s_{i_1}\cdots s_{i_k}y$.
    Then define
    $$ \lambda(x,v) = (y,\sigma^{i_k}\cdots\sigma^{i_1}v) . $$
    Similarly define $\mu\colon\abs X\to\abs X$, which for $(x,v)$ maps
    $$ \mu(x,v) = (d_{i_k}\cdots d_{i_1}x,u) $$
    where $v=\delta^{i_1}\cdots\delta^{i_k}u$.

    Clearly $\lambda,\mu$ both map points to equivalent points.
    Then $\lambda\circ\mu$ maps a point to an equivalent non-degenerate point.

    Uniqueness is because it can be verified that if $p\sim p'$ then $\lambda\mu(p)=\lambda\mu(p')$.
    \qed

\eproof

An important fact is the following.

\bthrm

    The geometric realization of a simplicial set $X$ is a CW complex.

\ethrm

\bproof

    Take the $n$-cells of $\abs X$ to be the images of the non-degenerate $n$-simplices of $X$.
    That is, simplices $(x,\abs{\Delta^n})$ for $x\in X_n$ non-degenerate.
    By the previous lemma, these cover $\abs X$, and it can be verified that these make $\abs X$ into a CW complex.
    \qed

\eproof

\bthrm

    There is an adjunction $\abs\bul\colon\sSet\tof\Top\colon\S$ between the geometric realization and the singular set functors.

\ethrm

\bproof

    For a simplicial set $X$ and a topological space $Y$, we define maps
    $$ \Phi\colon\hom_\Top(\abs X,Y) \tof \hom_\sSet(X,\S(Y))\colon\Psi $$
    as follows.
    For a continuous $f\colon\abs X\to Y$, restricting $f$ to a non-degenerate simplex $(x,\abs{\Delta^n})$ restricts to a continuous map $\Phi(f)_n\colon\abs{\Delta^n}\to Y$.

    Conversely for a map $g\colon X\to\S(Y)$, for each $n$-simplex $x\in X$ we have a continuous function $\rho_x\colon\abs{\Delta^n}\to Y$.
    So define $\Psi(g)\colon\abs X\to Y$ to act on the simplex $(x,\abs{\Delta^n})\in\abs X$ by applying $\rho_x$ to $\abs{\Delta^n}$.
    \qed

\eproof

